\section{\textcolor{blue}{Conclusiones}}


\begin{frame}{\textcolor{blue}{Conclusiones}}
  \begin{itemize}
      \item \textcolor{green}{Logros}
      \begin{enumerate}[(1)]
    \item  Se dio un ejemplo de que no toda función vectorial $ f: [a,b] \longrightarrow X  $ creciente es débil de variación acotada.
    \item  Se probo que para toda función vectorial $ f: [a,b] \longrightarrow X  $ creciente sea débil de variación acotada, es equivalente a que el cono sea de normal.
   % \item  Se probo que un cono sea sólido es condición suficiente, pero no necesaria para que $X$ cumpla la descomposición de Jordan dando un ejemplo.
    \end{enumerate}
    \item \textcolor{purple}{Perspectivas}
    \begin{itemize}
        \item[(1)] Ver la relación que el cono sea normal a que toda función vectorial sea integrable. 
        \item[(2)]  Ver si es equivalente que el cono extendible a que toda función vectorial sea  de variación acotada. 
        %\item[(3)] Analizar más a profundidad el espacio $ BV([a,b],X) $ como ENO  y aplicar los teoremas de punto fijo que tenemos.       
    \end{itemize}
  \end{itemize}
    
     \end{frame}

